\documentclass[10pt,a4paper]{amsart}
\usepackage[margin=19mm]{geometry}
\usepackage{amsmath,amssymb,mathtools}
\usepackage[scr=rsfs,cal=euler]{mathalfa}
\usepackage{xfrac}
\usepackage[unicode,hidelinks,bookmarks=false]{hyperref}
\newcommand{\ie}{i.e.\ }
\newcommand{\cf}{cf.\ }
\newcommand{\resp}{respectively\ }
\newcommand{\Subsection}{\subsection}
\providecommand{\nobreakdash}{\nobreak\hbox{-}}
\setlength{\emergencystretch}{2em}
\multlinegap0pt
\usepackage{bm}
\usepackage{bbm}
\usepackage{dsfont}
\usepackage{xspace}
\usepackage{cancel}
\usepackage{mathtools}
\usepackage{amsthm}
\makeatletter
\@ifundefined{theorem}{\newtheorem{theorem}{Theorem}}{}
\makeatother
\newcommand{\Z}{\mathbb{Z}}
\title[What if we decompose a simple tone? The Chinese remainder th]{What if we decompose a simple tone? The Chinese remainder theorem and structured Levi graphs in music}
\author{Pawe{\l} Nurowski}
\date{Source: arXiv:2605.25649v1 (CC BY 4.0)}
\begin{document}
\maketitle

\noindent\emph{Source and licence.} What if we decompose a simple tone? The Chinese remainder theorem and structured Levi graphs in music. Version arXiv:2605.25649v1 (CC BY 4.0), distributed under the Creative Commons Attribution 4.0 International licence, as recorded in the arXiv licence metadata for this exact version and shown on the abstract page https://arxiv.org/abs/2605.25649v1. Full rights notice: licenses/arxiv-2605.25649-CC-BY-4.0.txt.\par\smallskip

\noindent\textit{Editorial scope.} Counted unit: the Theorem whose source label is \texttt{None} in the article (source0.tex lines 333--339, source label \texttt{None}), delivered together with the article's own complete original proof of it (source0.tex lines 341--378, one \texttt{proof} environment of 38 lines). No mathematics is written, repaired or re-derived: statement and proof are the article's own text line for line. Required context retained uncounted: none. Every definition, display and label of those lines is retained unchanged. Omitted from this package and not redistributed: 1--332, 340--340, 379--2461. Nothing inside a retained excerpt is omitted or altered: every retained body is delivered exactly as the article's lines are written, so no omission receipt of any kind accompanies this package. Citations: the retained excerpts carry no citation command at all, so no bibliography entry of the article is transcribed and no reference is redistributed. Reference adaptation: none was needed and none was made; every reference inside the delivered unit resolves inside it, so no unresolved reference or citation is produced. Printed slips kept literal: no printed word, symbol, hypothesis or number of the counted statement or of its proof was altered. Layout and class: the article's own class is \texttt{amsart} with packages including fontenc, inputenc, babel, lmodern, placeins, float, geometry, hyperref, booktabs, enumitem, url, subcaption, xcolor, colortbl; the wrapper is the lane's pinned \texttt{amsart} editorial wrapper at 10pt on A4, which restates the theorem-like environments the retained text needs and adds the article's own macro definitions that the retained text needs (\texttt{\textbackslash Z}), with extra wrapper packages (bm, bbm, dsfont, xspace, cancel, mathtools). The delivered numbering is the unit's own. Assets: no retained span contains a figure environment, a graphics-inclusion command, a PDF-page inclusion, a table or a TikZ picture; no image file is copied into this package, the package declares no asset dependency and no image file is redistributed with it. Licence: version arXiv:2605.25649v1 of this article is distributed under the Creative Commons Attribution 4.0 International licence, as recorded in the arXiv licence metadata for this exact version and shown on the abstract page https://arxiv.org/abs/2605.25649v1. A full rights notice is delivered at licenses/arxiv-2605.25649-CC-BY-4.0.txt. Characters: every retained excerpt is pure ASCII, so the delivered engine, XeLaTeX with its default Latin Modern text font, reports no missing character.
\begin{theorem}[Note-Induced Isomorphisms in 12-TET]
In $\Z_{12}$, the condition of non-degeneracy yields exactly four harmonic systems with isomorphic Levi graphs: $(4,3)$, $(8,3)$, $(9,4)$, and $(9,8)$. Every graph isomorphism between any pair of these four systems can be explicitly induced by an affine note bijection $f(x) \equiv ax + b \pmod{12}$. Specifically:
\begin{enumerate}
    \item The systems $(4,3)$, $(8,3)$, and $(9,4)$ form an \textbf{orientation-preserving family}. They are mutually interconnected by affine bijections that strictly preserve chord polarities (Major maps to Major). 
    \item The system $(9,8)$ acts as the \textbf{orientation-reversing mirror}. Any affine note bijection connecting the orientation-preserving family to $(9,8)$ necessarily reverses chord polarities (Major maps to Minor).
\end{enumerate}
\end{theorem}

\begin{proof}
To prove that a Levi graph isomorphism is induced by notes, we must demonstrate an affine bijection $f(x) \equiv ax+b \pmod{12}$ that maps the pitch classes of the chords in the source system directly to the pitch classes of the chords in the target system. Let us take the classical system $(4,3)$ as our base. Its Major chord is $M_k^{(4,3)} = \{k, k+4, k+7\}$ and its Minor chord is $m_k^{(4,3)} = \{k, k+3, k+7\}$.

\textbf{1. Passage from (4,3) to (8,3):}
Let $f(x) = 5x$. Applying this element-wise to the notes of the $(4,3)$ chords yields:
\begin{align*}
    f(M_k^{(4,3)}) &= \{5k, 5k+20, 5k+35\} \equiv \{5k, 5k+8, 5k+11\} \pmod{12}. \\
    f(m_k^{(4,3)}) &= \{5k, 5k+15, 5k+35\} \equiv \{5k, 5k+3, 5k+11\} \pmod{12}.
\end{align*}
In the $(8,3)$ system, the generator is $q=11$, so its defining chords are $M_j^{(8,3)} = \{j, j+8, j+11\}$ and $m_j^{(8,3)} = \{j, j+3, j+11\}$. 
Substituting $j = 5k$, we see that $f(M_k^{(4,3)}) = M_{5k}^{(8,3)}$ and $f(m_k^{(4,3)}) = m_{5k}^{(8,3)}$. The map $f(x)=5x$ perfectly induces the Levi graph isomorphism while preserving chord polarity (orientation-preserving).

\textbf{2. Passage from (4,3) to (9,4):}
Let $f(x) = 5x+1$. Applying this to the $(4,3)$ chords yields:
\begin{align*}
    f(M_k^{(4,3)}) &= \{5k+1, 5k+21, 5k+36\} \equiv \{5k+1, 5k+9, 5k\} \pmod{12}. \\
    f(m_k^{(4,3)}) &= \{5k+1, 5k+16, 5k+36\} \equiv \{5k+1, 5k+4, 5k\} \pmod{12}.
\end{align*}
In the $(9,4)$ system, $q=1$, so $M_j^{(9,4)} = \{j, j+9, j+1\}$ and $m_j^{(9,4)} = \{j, j+4, j+1\}$. 
Letting $j = 5k$, we find $f(M_k^{(4,3)}) = M_{5k}^{(9,4)}$ and $f(m_k^{(4,3)}) = m_{5k}^{(9,4)}$. The map $f(x)=5x+1$ induces an orientation-preserving isomorphism.

\textbf{3. Passage from (4,3) to (9,8):}
Let $f(x) = 11x$. 
\begin{align*}
    f(M_k^{(4,3)}) &= \{11k, 11k+44, 11k+77\} \equiv \{11k, 11k+8, 11k+5\} \pmod{12}. 
\end{align*}
In the $(9,8)$ system, $q=5, t=9, s=8$, so a Minor chord is $m_j^{(9,8)} = \{j, j+8, j+5\}$. 
Thus, $f(M_k^{(4,3)}) = m_{11k}^{(9,8)}$. A Major chord has been mapped to a Minor chord. The map $f(x)=11x$ induces an orientation-reversing isomorphism (Negative Harmony).

\textbf{4. Passage from (8,3) to (9,8):}
To map from $(8,3)$ directly to $(9,8)$, we compose the inverse of the $(4,3)\to(8,3)$ map with the $(4,3)\to(9,8)$ map. The required multiplier is $11 \times 5^{-1} \equiv 11 \times 5 \equiv 55 \equiv 7 \pmod{12}$. Let $f(x) = 7x$.
\begin{align*}
    f(M_k^{(8,3)}) &= 7 \times \{k, k+8, k+11\} = \{7k, 7k+56, 7k+77\} \equiv \{7k, 7k+8, 7k+5\} \pmod{12}.
\end{align*}
This is precisely $m_{7k}^{(9,8)}$. Thus, $f(x)=7x$ induces an orientation-reversing isomorphism between $(8,3)$ and $(9,8)$.

By explicitly constructing the affine note bijections covering all cases, the theorem is proven.
\end{proof}

\end{document}
